\section{Discussion}\label{discussion}

\subsection{Three converging lines of evidence}\label{three-converging-lines-of-evidence}

The causal, behavioral, and mechanistic results converge on a single conclusion: the moral probe directions identified in \citet{reblitzrichardson2026geometry} are genuine features of the model's computation, not artifacts of the probing procedure.

Direction ablation establishes \emph{necessity}: the model needs these directions to generate foundation-specific content.
Steering injection establishes \emph{sufficiency}: adding these directions shifts model behavior toward the target foundation, with a dose--response curve characteristic of real feature manipulation.
Behavioral grounding establishes \emph{functional accessibility}: the directions predict which foundation a novel text activates, at accuracy levels far above chance.
SAE analysis establishes \emph{mechanistic correspondence}: unsupervised feature discovery partially recovers the same subspace that supervised probing identifies.

No single line of evidence would be conclusive.
Ablation could reflect a generic disruption effect; injection could be noise-driven; behavioral accuracy could reflect confounds; SAE overlap could be coincidental.
But taken together, the convergence across independent methods with different assumptions and failure modes provides strong evidence for the reality of the moral direction structure.

\subsection{The care saturation phenomenon}\label{the-care-saturation-phenomenon}

The MFV results reveal a property of real-world moral stimuli that laboratory-controlled probing datasets obscure: witnessing moral transgressions of \emph{any} kind activates the care/harm representation.
This is not a deficiency of the probing directions but a genuine feature of moral cognition as encoded by the model.
Moral psychology has long recognized that harm is the most accessible moral dimension cross-culturally \citep{graham2013mft}, and the model's representations reflect this asymmetry.

This has practical implications for steering applications: injecting a non-care foundation direction into a context that already contains harm content will contend with strong care co-activation.
Effective foundation-specific steering may require \emph{suppressing} the care component while boosting the target foundation---a two-direction intervention rather than a single injection.

\subsection{Layer-dependent causal roles}\label{layer-dependent-causal-roles}

Ablation and injection reveal complementary layer-dependent patterns:

\begin{itemize}
\tightlist
\item
  \textbf{Ablation specificity increases with depth} (strongest at layer 12), indicating that later layers rely more heavily on foundation-specific information for generation.
\item
  \textbf{Injection at layer 12 tolerates higher \(\alpha\)} before saturating, consistent with later-layer representations being more robust to perturbation.
\end{itemize}

These patterns suggest a processing hierarchy: earlier layers encode moral content in a form that is more malleable but less causally connected to output, while later layers encode it in a form that is more rigid but directly drives generation.

\subsection{Toward a steering fitness function}\label{toward-a-steering-fitness-function}

The results suggest a concrete fitness function for evaluating moral directions as steering targets:

\begin{enumerate}
\def\labelenumi{\arabic{enumi}.}
\tightlist
\item
  \textbf{Ablation specificity} \(< -0.5\) at the target layer (direction is causally load-bearing).
\item
  \textbf{Injection specificity} \(> 0\) at \(\alpha = 1\) (direction can produce specific behavioral shifts).
\item
  \textbf{Behavioral accuracy} \(> 50\%\) on held-out data (direction is functionally discriminative).
\item
  \textbf{SAE subspace overlap} \(> 2\times\) random baseline (direction corresponds to native features).
\end{enumerate}

Sanctity, loyalty, and care directions pass all four criteria.
Fairness and liberty pass criteria 1--3 but show weaker SAE overlap, suggesting their representations may be more distributed.
Authority passes all criteria at moderate levels.

\subsection{Limitations}\label{limitations}

\textbf{SAE training scale.} The 2M-token SAE achieves only 71.5\% variance explained.
Production-scale SAE training (100M+ tokens) would likely increase both feature quality and moral direction overlap.
The current results should be interpreted as a lower bound.

\textbf{Single model.} All experiments use OLMo-2 1B.
Generalization to larger models and different architectures is untested.

\textbf{Causal evaluation set size.} The 48-prompt evaluation set (8 per foundation) limits statistical power for per-foundation claims.
The direction of all effects is consistent, but effect sizes should be interpreted with appropriate uncertainty.

\textbf{Steering injection confounds.} At high \(\alpha\), injection degrades model performance broadly.
The positive specificity at high \(\alpha\) reflects \emph{relative} preservation of on-target content, not \emph{absolute} improvement.
The therapeutically relevant regime is \(\alpha = 1\)--\(2\).
